% ============================================================
% Capítulo 6 — Experimentos y resultados
% ============================================================

\chapter{Experimentos y resultados}
\label{cap:resultados}

Este capítulo presenta los resultados en el orden de la estrategia del
Capítulo~\ref{cap:planteamiento}: primero el post-proceso clásico de la
baseline (E1), después la comparación de pérdidas (E3), con especial
atención a Topograph, y por último el ajuste de hiperparámetros. Salvo
que se diga otra cosa, las cifras son de val~$B$, con el protocolo del
Capítulo~\ref{cap:metodologia}. Las cifras de val~$A$ se muestran porque
son las que guían las decisiones.

Conviene tener presente desde el principio la escala del ruido.
Reentrenar la misma configuración cambia el PQ en $B$ entre $0{,}01$ y
$0{,}02$ (Sección~\ref{sec:evaluacion}), y cambiar la semilla puede
moverlo más: la baseline pasa de $0{,}311$ con la semilla 42 a $0{,}287$
con la 43. Las diferencias entre pérdidas deben leerse con esa escala.

% ------------------------------------------------------------
\section{Post-proceso clásico (E1)}
\label{sec:res-post}
% ------------------------------------------------------------

La Tabla~\ref{tab:res-post} recoge las variantes del post-proceso sobre
la misma baseline (semilla 42), sin reentrenar.

\begin{table}[ht]
\centering
\small
\begin{tabular}{@{}lccccccc@{}}
\toprule
Post-proceso & $t$ & $a$ & PQ & SQ & RQ & TP & FP / FN\\
\midrule
Sin post-proceso & 0,50 & 0 & 0,311 & 0,638 & 0,488 & 495 & 665 / 374\\
B1 (rejilla) & 0,65 & 0 & 0,314 & 0,647 & 0,485 & 494 & 674 / 375\\
B1 (Optuna) & 0,66 & 0 & 0,313 & 0,646 & 0,484 & 494 & 678 / 375\\
B2 (rejilla) & 0,65 & 25 & \textbf{0,375} & 0,649 & 0,577 & 473 & 297 / 396\\
B2 (Optuna) & 0,66 & 29 & 0,373 & 0,650 & 0,574 & 458 & 269 / 411\\
\bottomrule
\end{tabular}
\caption{Post-proceso de la baseline (semilla 42) en val~$B$. $t$ es el
umbral y $a$ el área mínima en píxeles a $512\times512$, ambos ajustados
en val~$A$.}
\label{tab:res-post}
\end{table}

Se observan tres cosas:
\begin{itemize}
  \item \emph{El umbral por sí solo no aporta nada.} Subirlo de $0{,}5$
  a $0{,}65$ deja el PQ igual. Incluso aumenta algo los falsos
  positivos, probablemente porque al subir el umbral algunas
  componentes se parten en varias.
  \item \emph{El área mínima lo aporta casi todo.} B2 sube el PQ en
  $0{,}06$, de $0{,}311$ a $0{,}373$. Los falsos positivos bajan de 665
  a 269, a cambio de perder 37 aciertos. Es decir, la mayor parte de los
  falsos positivos de la baseline son componentes pequeñas.
  \item \emph{La mejora está en la detección, no en la delimitación.}
  SQ apenas cambia (de $0{,}64$ a $0{,}65$), mientras que RQ pasa de
  $0{,}49$ a $0{,}57$.
\end{itemize}
Con B2 la baseline supera el umbral de $0{,}35$ que los organizadores
consideran un resultado de gran valor (Sección~\ref{sec:competicion}).
Rejilla y Optuna llegan a valores muy parecidos; en lo que sigue se
usa B2 con Optuna para todas las ejecuciones.

% ------------------------------------------------------------
\section{Comparación de pérdidas (E3)}
\label{sec:res-perdidas}
% ------------------------------------------------------------

La Tabla~\ref{tab:res-perdidas} compara todas las pérdidas de la
Tabla~\ref{tab:perdidas}, sin post-proceso y con B2.

\begin{table}[ht]
\centering
\small
\begin{tabular}{@{}lccccc@{}}
\toprule
 & & \multicolumn{2}{c}{Sin post-proceso} & \multicolumn{2}{c}{B2 (Optuna)}\\
\cmidrule(lr){3-4}\cmidrule(lr){5-6}
Pérdida & Época & PQ $A$ & PQ $B$ & PQ $A$ & PQ $B$\\
\midrule
Baseline (s42) & 8 & 0,297 & 0,311 & \textbf{0,364} & \textbf{0,373}\\
Baseline (s43) & 8 & 0,276 & 0,287 & 0,348 & 0,357\\
IoU & 15 & 0,279 & 0,285 & 0,351 & 0,351\\
clDice & 13 & 0,307 & 0,313 & 0,351 & 0,354\\
IoU+clDice & 14 & 0,308 & 0,322 & 0,345 & 0,363\\
clDice sin Dice & 1 & 0,012 & 0,009 & 0,062 & 0,036\\
\midrule
Topograph (s42) & 17 & \textbf{0,322} & \textbf{0,327} & 0,354 & 0,354\\
Topograph (s43) & 12 & 0,298 & 0,322 & 0,336 & 0,355\\
IoU+Topo & 16 & 0,315 & 0,315 & 0,346 & 0,346\\
clDice+Topo & 18 & 0,307 & 0,317 & 0,344 & 0,348\\
IoU+clDice+Topo & 18 & 0,297 & 0,318 & 0,333 & 0,336\\
clDice sin Dice+Topo & 6 & 0,025 & 0,014 & 0,080 & 0,045\\
\bottomrule
\end{tabular}
\caption{PQ de cada pérdida en val~$A$ y val~$B$, sin post-proceso
(umbral $0{,}5$, sin filtro de área) y con B2 ajustado en val~$A$ para
cada ejecución. «Época» es la elegida en val~$A$. Semilla 42 salvo que
se indique.}
\label{tab:res-perdidas}
\end{table}

\paragraph{clDice sin Dice colapsa.} Con solo BCE y clDice, el mejor
modelo en $A$ es el de la primera época (la sexta con Topograph), y su
PQ no llega a $0{,}05$ ni con post-proceso. Es coherente con lo visto en
la Sección~\ref{sec:dice-cldice}: clDice compara esqueletos y no
controla la extensión de la máscara, así que necesita un término de
solapamiento que lo haga. En lo que sigue se dejan fuera estas dos
ejecuciones.

\paragraph{Sin post-proceso, las pérdidas con sesgo topológico ganan algo.}
Todas las variantes con clDice o con Topograph superan a la baseline de
la semilla 42, entre $+0{,}002$ (clDice) y $+0{,}016$ (Topograph). IoU,
que no tiene sesgo topológico, queda por debajo. La ventaja más clara es
la de Topograph, porque se repite con las dos semillas: $+0{,}016$ con
la 42 ($0{,}327$ frente a $0{,}311$) y $+0{,}035$ con la 43 ($0{,}322$
frente a $0{,}287$). Como se verá en la
Sección~\ref{sec:res-bootstrap}, solo la primera de estas diferencias
es estadísticamente distinguible de cero.

\paragraph{Con B2, la ventaja desaparece.} Tras el post-proceso, ninguna
pérdida supera a la baseline. Topograph queda en $0{,}354$ y $0{,}355$
con las dos semillas, frente a $0{,}373$ y $0{,}357$ de la baseline. La
mejor alternativa es IoU+clDice ($0{,}363$), dentro del rango que separa
las dos semillas de la baseline. Si se aplica el protocolo y se elige en
$A$, el modelo elegido es la baseline con B2.

\paragraph{Añadir Topograph a otras pérdidas no ayuda.} Sin
post-proceso, Topograph mejora a IoU ($+0{,}030$) pero apenas cambia
clDice ($+0{,}004$) e IoU+clDice ($-0{,}004$). Con B2, todas las
variantes con Topograph quedan por debajo de su versión sin él.

% ------------------------------------------------------------
\section{Qué hace Topograph}
\label{sec:res-topograph}
% ------------------------------------------------------------

Para entender por qué la ventaja de Topograph desaparece con el
post-proceso, la Tabla~\ref{tab:res-topo} descompone las detecciones.

\begin{table}[ht]
\centering
\small
\begin{tabular}{@{}llcccc@{}}
\toprule
Pérdida & Post & PQ & TP & FP & FN\\
\midrule
Baseline (s42) & sin & 0,311 & 495 & 665 & 374\\
Topograph (s42) & sin & 0,327 & 461 & 481 & 408\\
Baseline (s43) & sin & 0,287 & 503 & 899 & 366\\
Topograph (s43) & sin & 0,322 & 422 & 402 & 447\\
\midrule
Baseline (s42) & B2 & 0,373 & 458 & 269 & 411\\
Topograph (s42) & B2 & 0,354 & 457 & 342 & 412\\
Baseline (s43) & B2 & 0,357 & 431 & 279 & 438\\
Topograph (s43) & B2 & 0,355 & 412 & 220 & 457\\
\bottomrule
\end{tabular}
\caption{Detecciones en val~$B$ de la baseline y de Topograph, con las
dos semillas, sin post-proceso y con B2.}
\label{tab:res-topo}
\end{table}

Sin post-proceso, Topograph reduce mucho los falsos positivos: un
28\,\% con la semilla 42 (de 665 a 481) y un 55\,\% con la 43 (de 899 a
402). A cambio pierde aciertos, un 7\,\% y un 16\,\%. El balance es
positivo porque los falsos positivos eran el error dominante.

Pero eso es justo lo que hace el filtro de área. Tras B2, la baseline
tiene menos falsos positivos que Topograph sin post-proceso (269 y 279,
frente a 481 y 402) con un número parecido de aciertos, y los dos
modelos quedan en valores de PQ similares. La lectura topológica de la
Sección~\ref{sec:lectura-topologica} sugiere una explicación: los falsos
positivos pequeños son componentes de $H_0$ que no están en la
referencia. Topograph penaliza en el entrenamiento las regiones críticas
que crean esas componentes, y el filtro de área las elimina después.
Ambos corrigen el mismo tipo de error, así que sus mejoras no se suman.

La Tabla~\ref{tab:res-asignaciones} descompone los errores en los
tipos de la Sección~\ref{sec:diagnostico}: filamentos fragmentados
(cubiertos por varias predicciones), fusiones (una predicción que cubre
varios filamentos), filamentos sin ninguna predicción que los toque
(perdidos) y predicciones que no tocan ningún filamento (espurias).

\begin{table}[ht]
\centering
\small
\begin{tabular}{@{}llcccc@{}}
\toprule
Pérdida & Post & Fragm.\ (\%) & Fusiones (\%) & Perdidos & Espurias\\
\midrule
Baseline (s42) & sin & 10,6 & 1,0 & 140 & 319\\
Baseline (s43) & sin & 12,8 & 1,0 & 100 & 482\\
clDice & sin & 8,4 & 1,5 & 122 & 278\\
IoU+clDice & sin & 6,6 & 1,4 & 124 & 261\\
Topograph (s42) & sin & 9,0 & 1,4 & 183 & 169\\
Topograph (s43) & sin & 7,8 & 1,6 & 224 & 110\\
\midrule
Baseline (s42) & B2 & 4,9 & 1,7 & 250 & 68\\
Baseline (s43) & B2 & 4,3 & 1,8 & 275 & 88\\
clDice & B2 & 4,5 & 1,9 & 213 & 119\\
IoU+clDice & B2 & 3,9 & 1,9 & 201 & 110\\
Topograph (s42) & B2 & 4,7 & 1,5 & 217 & 108\\
Topograph (s43) & B2 & 3,0 & 2,2 & 300 & 49\\
\bottomrule
\end{tabular}
\caption{Tipos de error en val~$B$ (definiciones oficiales, IoU $>0$).
Fragmentados, en \% de los 869 filamentos de referencia; fusiones, en
\% de las predicciones.}
\label{tab:res-asignaciones}
\end{table}

Sin post-proceso, Topograph divide por dos o por tres las predicciones
espurias (de 319 y 482 a 169 y 110), pero a cambio deja más filamentos
sin tocar (de 140 y 100 a 183 y 224). También reduce algo la
fragmentación, aunque menos que clDice: con IoU+clDice los
fragmentados bajan del 10,6\,\% al 6,6\,\%, lo que concuerda con que
clDice premia la continuidad de las estructuras alargadas. Las fusiones
son raras en todos los casos (1--2\,\%) y aumentan ligeramente con las
pérdidas topológicas.

Con B2, la fragmentación queda entre el 3 y el 5\,\% en todos los
modelos, y las diferencias entre pérdidas se reducen a la proporción
entre perdidos y espurias. Es decir, el filtro de área elimina también
los trozos pequeños de los filamentos fragmentados, de modo que la
ventaja de clDice y de Topograph en fragmentación deja de notarse.
Destaca IoU+clDice, que con B2 es el modelo con menos filamentos
perdidos (201) y más aciertos (474).

% ------------------------------------------------------------
\section{Incertidumbre de las comparaciones}
\label{sec:res-bootstrap}
% ------------------------------------------------------------

La Tabla~\ref{tab:res-bootstrap} da la diferencia de PQ de cada
ejecución con la baseline (semilla 42) y su intervalo del 95\,\% por
\emph{bootstrap} pareado por días.

\begin{table}[ht]
\centering
\small
\begin{tabular}{@{}lcccc@{}}
\toprule
 & \multicolumn{2}{c}{Sin post-proceso} & \multicolumn{2}{c}{B2 (Optuna)}\\
\cmidrule(lr){2-3}\cmidrule(lr){4-5}
Pérdida & $\Delta$PQ & IC 95\,\% & $\Delta$PQ & IC 95\,\%\\
\midrule
\InputIfFileExists{tablas/bootstrap}{}{\multicolumn{5}{c}{\textbf{[ejecutar scripts/resultados.py]}}\\}
\bottomrule
\end{tabular}
\caption{Diferencia de PQ en val~$B$ con la baseline (semilla 42) e
intervalo del 95\,\% (\emph{bootstrap} pareado por días, 2.000
remuestreos). $^*$: el intervalo no contiene el 0.}
\label{tab:res-bootstrap}
\end{table}

Sin post-proceso, solo Topograph con la semilla 42 mejora a la
baseline de forma significativa ($+0{,}016$, IC $[+0{,}001,\,+0{,}032]$),
e IoU es significativamente peor. Con B2 ocurre lo contrario: casi todas
las pérdidas quedan significativamente por debajo de la baseline.

Pero estos intervalos deben leerse con cuidado. La baseline con la
semilla 43 es significativamente peor que con la 42 sin post-proceso
($-0{,}024$, IC $[-0{,}039,\,-0{,}009]$), y casi lo es con B2
($-0{,}016$). Es la misma configuración: lo único que cambia es la
semilla. El \emph{bootstrap} por días mide cuánto depende el PQ de qué
días forman val~$B$, pero no la variabilidad del propio entrenamiento.
Con una sola ejecución por pérdida, una diferencia «significativa»
puede deberse a que la ejecución de referencia salió especialmente bien
o mal.

Tomando como referencia las dos semillas, la conclusión es más
modesta. Sin post-proceso, Topograph supera a la baseline con las dos
semillas, lo que apunta a una mejora real. Con B2, la baseline queda en
$0{,}357$--$0{,}373$ y Topograph en $0{,}354$--$0{,}355$: Topograph no
mejora, y la diferencia con la baseline es del orden de la que separa a
las dos semillas. Confirmarlo exigiría varias semillas por pérdida.

% ------------------------------------------------------------
\section{Barbas}
\label{sec:res-barbas}
% ------------------------------------------------------------

PQ no distingue si un acierto recoge las barbas. Antes de comparar
modelos conviene saber cuánto pesan las barbas en los datos y cuánto se
pierde solo por trabajar a $512\times512$.

\paragraph{Las barbas en val~$B$.} De los 869 filamentos anotados, el
41,4\,\% tiene al menos una barba, con una media de $0{,}70$ barbas por
filamento. Son estructuras pequeñas: su esqueleto mide en total un
6\,\% de la longitud de las espinas, y el grosor mediano de un
filamento es de $11{,}4$ píxeles a resolución completa (menos de 3 a
$512\times512$). Esto explica por qué PQ apenas las ve: un acierto con o
sin barbas tiene casi el mismo IoU.

\paragraph{Techo a $512\times512$.} Si se reduce cada máscara anotada a
$512\times512$ y se vuelve a ampliar, como le ocurre a la predicción del
modelo, y se mide contra la original, la cobertura de barbas es
$0{,}997$, la precisión $0{,}908$ y el $F_1$ $0{,}950$; se conservan
todas las barbas y aparecen $0{,}075$ ramas inventadas por filamento.
Dos consecuencias:
\begin{itemize}
  \item \emph{La resolución no es el cuello de botella.} Lo que el
  modelo no recupera de las barbas no se debe a la reducción de la
  imagen, sino a lo que aprende.
  \item \emph{La precisión tiene un techo de $0{,}91$.} El reescalado
  suaviza el contorno y crea pequeñas ramas en el esqueleto, que la
  medida cuenta como inventadas. Las precisiones de los modelos deben
  compararse con ese techo, no con 1.
\end{itemize}

\paragraph{Modelos.} La Tabla~\ref{tab:res-barbas} muestra las medidas
de barbas de la baseline, Topograph y clDice.

\begin{table}[ht]
\centering
\small
\begin{tabular}{@{}llcccc@{}}
\toprule
Pérdida & Post & Cobertura & Precisión & $F_1$ & Ramas inventadas\\
\midrule
Techo a $512\times512$ & --- & 0,997 & 0,908 & 0,950 & 0,075\\
\midrule
\multicolumn{6}{c}{\textbf{[completar con \texttt{barbas\_runs.csv}]}}\\
\bottomrule
\end{tabular}
\caption{Medidas de barbas en val~$B$. El techo es la propia anotación
reducida a $512\times512$ y vuelta a ampliar.}
\label{tab:res-barbas}
\end{table}

\textbf{[completar: medido sobre los reentrenamientos de referencia
(intento 0 del ajuste, misma config que la Tabla~\ref{tab:res-perdidas});
decirlo en el texto. Si clDice o Topograph cubren más barbas que la
baseline, y a qué coste en precisión.]}

% ------------------------------------------------------------
\section{Ajuste de hiperparámetros}
\label{sec:res-hp}
% ------------------------------------------------------------

\textbf{[completar mañana con \texttt{ajuste\_hp.py --resumen}: mejor
intento de cada modelo (valores, PQ $A$ y $B$, sin post y con B2),
número de intentos y podados. Comparar la ganancia en $B$ con el ruido
de $0{,}01$--$0{,}02$. Recordar que el PQ en $A$ del mejor intento está
inflado por la selección; la cifra válida es la de $B$.]}

% ------------------------------------------------------------
\section{Resumen}
\label{sec:res-resumen}
% ------------------------------------------------------------

\begin{itemize}
  \item El post-proceso clásico es la mejora más grande y más barata:
  el área mínima sube el PQ de la baseline de $0{,}311$ a $0{,}373$,
  porque elimina componentes espurias pequeñas.
  \item Sin post-proceso, las pérdidas con sesgo topológico mejoran a la
  baseline. Topograph es la más consistente (mejora con las dos
  semillas) y la única con una diferencia significativa.
  \item Topograph mejora sobre todo porque reduce las predicciones
  espurias, a costa de perder algunos filamentos. clDice reduce más la
  fragmentación.
  \item El filtro de área corrige esos mismos errores, así que con B2
  la ventaja de ambas desaparece. Ninguna pérdida supera a la baseline
  con B2.
  \item clDice necesita un término de solapamiento; sin Dice, el
  entrenamiento colapsa.
  \item La variación entre semillas es del mismo orden que las
  diferencias entre pérdidas, y el \emph{bootstrap} por días no la
  recoge. \textbf{[ajustar tras el ajuste de hiperparámetros.]}
\end{itemize}
